\scanpage{64}
% archive foliation: 63
Faisons les calculs pour $J=\{0,1,\infty\}\subset\mathbb P^1_S$, on trouve facilement:
\[
\tag{34}\left\{\begin{array}{ll}
\sigma_0^t(z)=\dfrac{z-t}{z-1}&(\text{éch. }0\text{ et }t,\ 1\text{ et }\infty)\\
\sigma_1^t(z)=t/z&(\text{éch. }1\text{ et }t,\ 0\text{ et }\infty)\\
\sigma_\infty^t(z)=t\dfrac{z-1}{z-t}&(\text{éch. }\infty\text{ et }t,\ 0\text{ et }1)
\end{array}\right.
\]
\margin{NB $\sigma_\infty^t=\sigma_0^t\sigma_1^t$}
On trouve donc
\[
\tag{35}\left\{\begin{array}{lll}
\Delta_0^t:&\text{solutions de }z^2-2z+t=0,&\Delta_0^t=\{1\pm\sqrt{1-t}\}\\
\Delta_1^t:&\text{\rule{3cm}{0.4pt}}\quad z^2-t=0&\Delta_1^t=\{\pm\sqrt t\}\\
\Delta_\infty^t:&\text{\rule{3cm}{0.4pt}}\quad z^2-2tz+t=0&\Delta_\infty^t=\{t\pm\sqrt{t(1-t)}\}
\end{array}\right.
\]
et on doit avoir
\[
\tag{36}f_2^t:\quad\begin{array}{l}
\Delta_0^t\longrightarrow0\\
\Delta_1^t\longrightarrow1\\
\Delta_\infty^t\longrightarrow\infty
\end{array}\qquad\uline{\text{avec multiplicité }2}
\]
ce qui nous donne (puisque les zéros et pôles de $f_2^t$ de degré quatre sont connus)
\[
f_2^t(z)=c_t\left(\frac{z^2-2z+t}{z^2-2tz+t}\right)^2\qquad c_t\text{ section de }\underline O_S^*,
\]
dont la valeur s'obtient en écrivant que $\Delta_1^t\longrightarrow1$. Mais si $\zeta$ est une section de $\Delta_1^t$ i.e. $\zeta^2=t$, on a
\margin{$\begin{aligned}
\zeta^2\deleted{{}+t}-2\zeta+t&=2(t-\zeta)\\
\zeta^2-2t\zeta+t&=2t(1-\zeta)
\end{aligned}$}
\[
f_2^t(\zeta)=c_t\left(\frac{2(t-\zeta)}{2t(1-\zeta)}\right)^2
=c_t/t^2\ \deleted{\frac{\zeta^2-2t\zeta+t^2}{\zeta^2-2\zeta+1}}
=\left(\frac{t-\zeta}{1-\zeta}\right)^2
\]
Or
\[
\deleted{{}=\frac{c_t}{t^2}\frac{t^2+t-2t\zeta}{1+t-2\zeta}}\quad-
\quad\frac{t-\zeta}{1-\zeta}
=\frac{(t-\zeta)(1+\zeta)}{(1-\zeta)(1+\zeta)}
=\frac{-\zeta^2+(t-1)\zeta+t}{1-\zeta^2}
=\frac{(t-1)\zeta}{1-t}=-\zeta
\]
donc
\[
f_2^t(\zeta)=\frac{c_t}{t^2}\zeta^2=\frac{c_t t}{t^2}=\frac{c_t}{t}\qquad\text{d'où}\quad c_t=t
\]
